\documentclass[10pt,a4paper]{amsart}
\usepackage{amsmath,amssymb,amsthm,graphicx}
% For commutative diagrams we use Xy-pic
\usepackage[all]{xy}

% We use 2cell for 2-commutative diagrams.
\xyoption{2cell}
% Theorem environments.
%
\theoremstyle{plain}
\newtheorem{theorem}[subsection]{Theorem}
\newtheorem{proposition}[subsection]{Proposition}
\newtheorem{lemma}[subsection]{Lemma}

\theoremstyle{definition}
\newtheorem{definition}[subsection]{Definition}
\newtheorem{example}[subsection]{Example}
\newtheorem{exercise}[subsection]{Exercise}
\newtheorem{situation}[subsection]{Situation}

\theoremstyle{remark}
\newtheorem{remark}[subsection]{Remark}
\newtheorem{remarks}[subsection]{Remarks}

\numberwithin{equation}{subsection}

% Macros
%
\def\lim{\mathop{\mathrm{lim}}\nolimits}
\def\colim{\mathop{\mathrm{colim}}\nolimits}
\def\Spec{\mathop{\mathrm{Spec}}}
\def\Hom{\mathop{\mathrm{Hom}}\nolimits}
\def\Ext{\mathop{\mathrm{Ext}}\nolimits}
\def\SheafHom{\mathop{\mathcal{H}\!\mathit{om}}\nolimits}
\def\SheafExt{\mathop{\mathcal{E}\!\mathit{xt}}\nolimits}
\def\Sch{\mathit{Sch}}
\def\Mor{\mathop{\mathrm{Mor}}\nolimits}
\def\Ob{\mathop{\mathrm{Ob}}\nolimits}
\def\Sh{\mathop{\mathit{Sh}}\nolimits}
\def\NL{\mathop{N\!L}\nolimits}
\def\CH{\mathop{\mathrm{CH}}\nolimits}
\def\proetale{{pro\text{-}\acute{e}tale}}
\def\etale{{\acute{e}tale}}
\def\QCoh{\mathit{QCoh}}
\def\Ker{\mathop{\mathrm{Ker}}}
\def\Im{\mathop{\mathrm{Im}}}
\def\Coker{\mathop{\mathrm{Coker}}}
\def\Coim{\mathop{\mathrm{Coim}}}

% Boxtimes
%
\DeclareMathSymbol{\boxtimes}{\mathbin}{AMSa}{"02}

%
% Macros for moduli stacks/spaces
%
\def\QCohstack{\mathcal{QC}\!\mathit{oh}}
\def\Cohstack{\mathcal{C}\!\mathit{oh}}
\def\Spacesstack{\mathcal{S}\!\mathit{paces}}
\def\Quotfunctor{\mathrm{Quot}}
\def\Hilbfunctor{\mathrm{Hilb}}
\def\Curvesstack{\mathcal{C}\!\mathit{urves}}
\def\Polarizedstack{\mathcal{P}\!\mathit{olarized}}
\def\Complexesstack{\mathcal{C}\!\mathit{omplexes}}
% \Pic is the operator that assigns to X its picard group, usage \Pic(X)
% \Picardstack_{X/B} denotes the Picard stack of X over B
% \Picardfunctor_{X/B} denotes the Picard functor of X over B
\def\Pic{\mathop{\mathrm{Pic}}\nolimits}
\def\Picardstack{\mathcal{P}\!\mathit{ic}}
\def\Picardfunctor{\mathrm{Pic}}
\def\Deformationcategory{\mathcal{D}\!\mathit{ef}}
\usepackage{cleveref}
\crefname{originalproposition}{original proposition}{original propositions}
\Crefname{originalproposition}{Original proposition}{Original propositions}
\makeatletter
\newlabel{proposition-finite-higher-direct-image-zero}{{03QP}{source}{}{sourcepointer-1}{}}
\newlabel{proposition-finite-higher-direct-image-zero@cref}{{[originalproposition][0][]03QP}{[1][0][]source}{}{}{}}
\makeatother
\crefname{originallemma}{original lemma}{original lemmas}
\Crefname{originallemma}{Original lemma}{Original lemmas}
\makeatletter
\newlabel{sites-cohomology-lemma-Leray}{{0732}{source}{}{sourcepointer-2}{}}
\newlabel{sites-cohomology-lemma-Leray@cref}{{[originallemma][0][]0732}{[1][0][]source}{}{}{}}
\makeatother
\crefname{originallemma}{original lemma}{original lemmas}
\Crefname{originallemma}{Original lemma}{Original lemmas}
\makeatletter
\newlabel{lemma-cd-curve-over-field}{{0F0S}{source}{}{sourcepointer-3}{}}
\newlabel{lemma-cd-curve-over-field@cref}{{[originallemma][0][]0F0S}{[1][0][]source}{}{}{}}
\makeatother
\crefname{originallemma}{original lemma}{original lemmas}
\Crefname{originallemma}{Original lemma}{Original lemmas}
\makeatletter
\newlabel{varieties-lemma-noether-normalization-affine}{{0CBI}{source}{}{sourcepointer-4}{}}
\newlabel{varieties-lemma-noether-normalization-affine@cref}{{[originallemma][0][]0CBI}{[1][0][]source}{}{}{}}
\makeatother
\crefname{originaltheorem}{original theorem}{original theorems}
\Crefname{originaltheorem}{Original theorem}{Original theorems}
\makeatletter
\newlabel{theorem-higher-direct-images}{{03Q9}{source}{}{sourcepointer-5}{}}
\newlabel{theorem-higher-direct-images@cref}{{[originaltheorem][0][]03Q9}{[1][0][]source}{}{}{}}
\makeatother
\crefname{originallemma}{original lemma}{original lemmas}
\Crefname{originallemma}{Original lemma}{Original lemmas}
\makeatletter
\newlabel{lemma-strictly-henselian}{{0F0U}{source}{}{sourcepointer-6}{}}
\newlabel{lemma-strictly-henselian@cref}{{[originallemma][0][]0F0U}{[1][0][]source}{}{}{}}
\makeatother
\crefname{originallemma}{original lemma}{original lemmas}
\Crefname{originallemma}{Original lemma}{Original lemmas}
\makeatletter
\newlabel{lemma-cd-field-extension}{{0F0T}{source}{}{sourcepointer-7}{}}
\newlabel{lemma-cd-field-extension@cref}{{[originallemma][0][]0F0T}{[1][0][]source}{}{}{}}
\makeatother
\crefname{originallemma}{original lemma}{original lemmas}
\Crefname{originallemma}{Original lemma}{Original lemmas}
\makeatletter
\newlabel{lemma-cd-limit}{{0F0R}{source}{}{sourcepointer-8}{}}
\newlabel{lemma-cd-limit@cref}{{[originallemma][0][]0F0R}{[1][0][]source}{}{}{}}
\makeatother
\crefname{originallemma}{original lemma}{original lemmas}
\Crefname{originallemma}{Original lemma}{Original lemmas}
\makeatletter
\newlabel{varieties-lemma-dimension-locally-algebraic}{{0A21}{source}{}{sourcepointer-9}{}}
\newlabel{varieties-lemma-dimension-locally-algebraic@cref}{{[originallemma][0][]0A21}{[1][0][]source}{}{}{}}
\makeatother
\crefname{originallemma}{original lemma}{original lemmas}
\Crefname{originallemma}{Original lemma}{Original lemmas}
\makeatletter
\newlabel{lemma-support-in-subset}{{0F0N}{source}{}{sourcepointer-10}{}}
\newlabel{lemma-support-in-subset@cref}{{[originallemma][0][]0F0N}{[1][0][]source}{}{}{}}
\makeatother
\crefname{originalproposition}{original proposition}{original propositions}
\Crefname{originalproposition}{Original proposition}{Original propositions}
\makeatletter
\newlabel{proposition-closed-immersion-pushforward}{{04CA}{source}{}{sourcepointer-11}{}}
\newlabel{proposition-closed-immersion-pushforward@cref}{{[originalproposition][0][]04CA}{[1][0][]source}{}{}{}}
\makeatother
\crefname{originaltheorem}{original theorem}{original theorems}
\Crefname{originaltheorem}{Original theorem}{Original theorems}
\makeatletter
\newlabel{theorem-colimit}{{09YQ}{source}{}{sourcepointer-12}{}}
\newlabel{theorem-colimit@cref}{{[originaltheorem][0][]09YQ}{[1][0][]source}{}{}{}}
\makeatother
\crefname{originallemma}{original lemma}{original lemmas}
\Crefname{originallemma}{Original lemma}{Original lemmas}
\makeatletter
\newlabel{lemma-cohomological-dimension-proper}{{095U}{source}{}{sourcepointer-13}{}}
\newlabel{lemma-cohomological-dimension-proper@cref}{{[originallemma][0][]095U}{[1][0][]source}{}{}{}}
\makeatother
\title{Cohomological dimension of an affine scheme of finite type over a field}
\author{The Stacks Project Authors}
\date{stacks-project:etale-cohomology.tex:local-snapshot-20261011; GNU FDL 1.2 or later}
\begin{document}
\maketitle
\noindent\textit{Editorial note.} All delivered mathematical text is the author's own, copied from the pinned Stacks Project chapter file etale-cohomology.tex; no statement or proof step was rewritten, reordered or removed.\par\medskip
\noindent\textit{Editorial note.} The proof is the author's complete argument, including both interludes and the final diagram; the Xy-pic diagram is delivered as the author wrote it.\par\medskip
\noindent\textit{Editorial note.} Every labelled reference occurring inside the counted spans points to a chapter of the Project that is not part of this card's source tree. Each is carried in assembly.external\_references as a non-load-bearing navigation pointer whose source location is the invocation inside the counted span, so no load-bearing mathematics is replaced by a stub.\par\medskip
\noindent\textit{Editorial note.} The checked chapter list for same-name collisions is the set of stacks chapters staged in this repository: chow, cohomology, divisors, etale-cohomology, flat, more-algebra, more-morphisms, morphisms, spaces-more-morphisms, varieties together with this card's chapter etale-cohomology.\par\medskip
\noindent\textit{Editorial note.} The preamble is the chapter's own native preamble excerpt already used for this source: the Xy-pic setup (native preamble lines 14-18) and the native theorem declarations with the mathematical macros (native preamble lines 154-217).\par\medskip
\noindent\textit{Allowed standard prerequisites (editorial).}\begin{itemize}
\item Cohomological dimension cd(X) of a scheme and the vanishing of torsion-sheaf cohomology. Same chapter, the definition of the cohomological dimension as the supremum of the degrees in which torsion sheaves can have nonzero cohomology. Supplies the quantity bounded in the statement and the reduction of the proof to vanishing statements for torsion sheaves.
\item Vanishing of higher direct images along a finite morphism. Same chapter, the result that a finite morphism has vanishing higher direct images. Supplies the case dim(X) = 0, where the structure morphism is finite.
\item Noether normalization for affine schemes of finite type over a field. Varieties, the Noether normalization producing a finite morphism from an affine variety to an affine space. Supplies the reduction of the general case to a smooth situation of the same dimension.
\item Vanishing for curves and the curve case of the bound. Same chapter, the bound cd(X) at most 1 + cd(K) for a one-dimensional affine scheme of finite type over K. Supplies the base of the induction on the dimension.
\end{itemize}\medskip
\ifdefined\hypertarget\hypertarget{sourcepointer-1}{}\fi
\noindent\textit{Source reference.} Original proposition 03QP: sources/stacks-etale-cohomology/pinned-original-source.tex, source lines 8245--8245. Navigation pointer invoked at proof line 19775; the counted span contains the pointer only and the cited result is not reproduced inside it, so the pointer is registered as a non-load-bearing navigation stub in the established form for Stacks material. No bibliography citation is involved.\par\medskip
\ifdefined\hypertarget\hypertarget{sourcepointer-2}{}\fi
\noindent\textit{Source reference.} Original lemma 0732: sources/stacks-etale-cohomology/pinned-original-source.tex, source lines 19778--19778. Navigation pointer invoked at proof line 19778; the counted span contains the pointer only and the cited result is not reproduced inside it, so the pointer is registered as a non-load-bearing navigation stub in the established form for Stacks material. No bibliography citation is involved.\par\medskip
\ifdefined\hypertarget\hypertarget{sourcepointer-3}{}\fi
\noindent\textit{Source reference.} Original lemma 0F0S: sources/stacks-etale-cohomology/pinned-original-source.tex, source lines 19651--19651. Navigation pointer invoked at proof line 19782; the counted span contains the pointer only and the cited result is not reproduced inside it, so the pointer is registered as a non-load-bearing navigation stub in the established form for Stacks material. No bibliography citation is involved.\par\medskip
\ifdefined\hypertarget\hypertarget{sourcepointer-4}{}\fi
\noindent\textit{Source reference.} Original lemma 0CBI: sources/stacks-etale-cohomology/pinned-original-source.tex, source lines 19788--19788. Navigation pointer invoked at proof line 19788; the counted span contains the pointer only and the cited result is not reproduced inside it, so the pointer is registered as a non-load-bearing navigation stub in the established form for Stacks material. No bibliography citation is involved.\par\medskip
\ifdefined\hypertarget\hypertarget{sourcepointer-5}{}\fi
\noindent\textit{Source reference.} Original theorem 03Q9: sources/stacks-etale-cohomology/pinned-original-source.tex, source lines 8015--8015. Navigation pointer invoked at proof line 19804; the counted span contains the pointer only and the cited result is not reproduced inside it, so the pointer is registered as a non-load-bearing navigation stub in the established form for Stacks material. No bibliography citation is involved.\par\medskip
\ifdefined\hypertarget\hypertarget{sourcepointer-6}{}\fi
\noindent\textit{Source reference.} Original lemma 0F0U: sources/stacks-etale-cohomology/pinned-original-source.tex, source lines 19692--19692. Navigation pointer invoked at proof line 19815; the counted span contains the pointer only and the cited result is not reproduced inside it, so the pointer is registered as a non-load-bearing navigation stub in the established form for Stacks material. No bibliography citation is involved.\par\medskip
\ifdefined\hypertarget\hypertarget{sourcepointer-7}{}\fi
\noindent\textit{Source reference.} Original lemma 0F0T: sources/stacks-etale-cohomology/pinned-original-source.tex, source lines 19673--19673. Navigation pointer invoked at proof line 19826; the counted span contains the pointer only and the cited result is not reproduced inside it, so the pointer is registered as a non-load-bearing navigation stub in the established form for Stacks material. No bibliography citation is involved.\par\medskip
\ifdefined\hypertarget\hypertarget{sourcepointer-8}{}\fi
\noindent\textit{Source reference.} Original lemma 0F0R: sources/stacks-etale-cohomology/pinned-original-source.tex, source lines 19633--19633. Navigation pointer invoked at proof line 19827; the counted span contains the pointer only and the cited result is not reproduced inside it, so the pointer is registered as a non-load-bearing navigation stub in the established form for Stacks material. No bibliography citation is involved.\par\medskip
\ifdefined\hypertarget\hypertarget{sourcepointer-9}{}\fi
\noindent\textit{Source reference.} Original lemma 0A21: sources/stacks-etale-cohomology/pinned-original-source.tex, source lines 19835--19835. Navigation pointer invoked at proof line 19835; the counted span contains the pointer only and the cited result is not reproduced inside it, so the pointer is registered as a non-load-bearing navigation stub in the established form for Stacks material. No bibliography citation is involved.\par\medskip
\ifdefined\hypertarget\hypertarget{sourcepointer-10}{}\fi
\noindent\textit{Source reference.} Original lemma 0F0N: sources/stacks-etale-cohomology/pinned-original-source.tex, source lines 13091--13091. Navigation pointer invoked at proof line 19842; the counted span contains the pointer only and the cited result is not reproduced inside it, so the pointer is registered as a non-load-bearing navigation stub in the established form for Stacks material. No bibliography citation is involved.\par\medskip
\ifdefined\hypertarget\hypertarget{sourcepointer-11}{}\fi
\noindent\textit{Source reference.} Original proposition 04CA: sources/stacks-etale-cohomology/pinned-original-source.tex, source lines 6702--6702. Navigation pointer invoked at proof line 19845; the counted span contains the pointer only and the cited result is not reproduced inside it, so the pointer is registered as a non-load-bearing navigation stub in the established form for Stacks material. No bibliography citation is involved.\par\medskip
\ifdefined\hypertarget\hypertarget{sourcepointer-12}{}\fi
\noindent\textit{Source reference.} Original theorem 09YQ: sources/stacks-etale-cohomology/pinned-original-source.tex, source lines 7400--7400. Navigation pointer invoked at proof line 19850; the counted span contains the pointer only and the cited result is not reproduced inside it, so the pointer is registered as a non-load-bearing navigation stub in the established form for Stacks material. No bibliography citation is involved.\par\medskip
\ifdefined\hypertarget\hypertarget{sourcepointer-13}{}\fi
\noindent\textit{Source reference.} Original lemma 095U: sources/stacks-etale-cohomology/pinned-original-source.tex, source lines 18786--18786. Navigation pointer invoked at proof line 19873; the counted span contains the pointer only and the cited result is not reproduced inside it, so the pointer is registered as a non-load-bearing navigation stub in the established form for Stacks material. No bibliography citation is involved.\par\medskip
\begin{proposition}
\label{proposition-cd-affine}
Let $K$ be a field. Let $X$ be an affine scheme of finite type over $K$.
Then we have $\text{cd}(X) \leq \dim(X) + \text{cd}(K)$.
\end{proposition}

\begin{proof}
We will prove this by induction on $\dim(X)$.
Let $\mathcal{F}$ be an abelian torsion sheaf on $X_\etale$.

\medskip\noindent
The case $\dim(X) = 0$. In this case the structure morphism
$f : X \to \Spec(K)$ is finite. Hence we see that $R^if_*\mathcal{F} = 0$
for $i > 0$, see Proposition \ref{proposition-finite-higher-direct-image-zero}.
Thus $H^i_\etale(X, \mathcal{F}) = H^i_\etale(\Spec(K), f_*\mathcal{F})$
by the Leray spectral sequence for $f$
(Cohomology on Sites, Lemma \ref{sites-cohomology-lemma-Leray})
and the result is clear.

\medskip\noindent
The case $\dim(X) = 1$. This is Lemma \ref{lemma-cd-curve-over-field}.

\medskip\noindent
Assume $d = \dim(X) > 1$ and the proposition holds for finite type
affine schemes of dimension $< d$ over fields.
By Noether normalization, see for example
Varieties, Lemma \ref{varieties-lemma-noether-normalization-affine},
there exists a finite morphism $f : X \to \mathbf{A}^d_K$.
Recall that $R^if_*\mathcal{F} = 0$ for $i > 0$ by
Proposition \ref{proposition-finite-higher-direct-image-zero}.
By the Leray spectral sequence for $f$
(Cohomology on Sites, Lemma \ref{sites-cohomology-lemma-Leray})
we conclude that it suffices to prove the result for $\pi_*\mathcal{F}$
on $\mathbf{A}^d_K$.

\medskip\noindent
Interlude I. Let $j : X \to Y$ be an open immersion of smooth
$d$-dimensional varieties over $K$ (not necessarily affine)
whose complement is the support of an effective Cartier divisor $D$.
The sheaves $R^qj_*\mathcal{F}$ for $q > 0$ are supported on $D$.
We claim that $(R^qj_*\mathcal{F})_{\overline{y}} = 0$
for $a = \text{trdeg}_K(\kappa(y)) > d - q$. Namely, by
Theorem \ref{theorem-higher-direct-images} we have
$$
(R^qj_*\mathcal{F})_{\overline{y}} =
H^q(\Spec(\mathcal{O}_{Y, y}^{sh}) \times_Y X, \mathcal{F})
$$
Choose a local equation $f \in \mathfrak m_y \subset \mathcal{O}_{Y, y}$
for $D$. Then we have
$$
\Spec(\mathcal{O}_{Y, y}^{sh}) \times_Y X =
\Spec(\mathcal{O}_{Y, y}^{sh}[1/f])
$$
Using Lemma \ref{lemma-strictly-henselian} we get an embedding
$$
K(t_1, \ldots, t_a)^{sep}(x) =
K(t_1, \ldots, t_a)^{sep}[x]_{(x)}[1/x]
\longrightarrow
\mathcal{O}_{Y, y}^{sh}[1/f]
$$
Since the transcendence degree over $K$ of the fraction field of
$\mathcal{O}_{Y, y}^{sh}$ is $d$, we see that $\mathcal{O}_{Y, y}^{sh}[1/f]$
is a filtered colimit of $(d - a - 1)$-dimensional finite type algebras over
the field $K(t_1, \ldots, t_a)^{sep}(x)$ which itself has cohomological
dimension $1$ by Lemma \ref{lemma-cd-field-extension}. Thus by induction
hypothesis and Lemma \ref{lemma-cd-limit}
we obtain the desired vanishing.

\medskip\noindent
Interlude II. Let $Z$ be a smooth variety over $K$ of dimension $d - 1$.
Let $E_a \subset Z$ be the set of points $z \in Z$ with
$\text{trdeg}_K(\kappa(z)) \leq a$. Observe that $E_a$ is closed
under specialization, see
Varieties, Lemma \ref{varieties-lemma-dimension-locally-algebraic}.
Suppose that $\mathcal{G}$ is a torsion abelian sheaf on $Z$
whose support is contained in $E_a$. Then we claim that
$H^b_\etale(Z, \mathcal{G}) = 0$ for $b > a + \text{cd}(K)$.
Namely, we can write $\mathcal{G} = \colim \mathcal{G}_i$
with $\mathcal{G}_i$ a torsion abelian sheaf
supported on a closed subscheme $Z_i$ contained in $E_a$, see
Lemma \ref{lemma-support-in-subset}.
Then the induction hypothesis kicks in to imply
the desired vanishing for $\mathcal{G}_i$\footnote{Here
we first use Proposition \ref{proposition-closed-immersion-pushforward}
to write $\mathcal{G}_i$ as the pushforward of a sheaf on $Z_i$,
the induction hypothesis gives the vanishing for this sheaf on $Z_i$, and
the Leray spectral sequence for $Z_i \to Z$ gives the vanishing
for $\mathcal{G}_i$.}. Finally, we
conclude by Theorem \ref{theorem-colimit}.

\medskip\noindent
Consider the commutative diagram
$$
\xymatrix{
\mathbf{A}^d_K \ar[rd]_f \ar[rr]_-j & &
\mathbf{P}^1_K \times_K \mathbf{A}^{d - 1}_K \ar[ld]^g \\
& \mathbf{A}^{d - 1}_K
}
$$
Observe that $j$ is an open immersion of smooth $d$-dimensional
varieties whose complement is an effective Cartier divisor $D$. Thus
we may use the results obtained in interlude I.
We are going to study the relative Leray spectral sequence
$$
E_2^{p, q} = R^pg_*R^qj_*\mathcal{F} \Rightarrow R^{p + q}f_*\mathcal{F}
$$
Since $R^qj_*\mathcal{F}$ for $q > 0$ is supported on $D$ and since
$g|_D : D \to \mathbf{A}^{d - 1}_K$ is an isomorphism, we find
$R^pg_*R^qj_*\mathcal{F} = 0$ for $p > 0$ and $q > 0$. Moreover, we have
$R^qj_*\mathcal{F} = 0$ for $q > d$. On the other hand, $g$ is a
proper morphism of relative dimension $1$. Hence by
Lemma \ref{lemma-cohomological-dimension-proper}
we see that $R^pg_*j_*\mathcal{F} = 0$ for $p > 2$. Thus the $E_2$-page
of the spectral sequence looks like this
$$
\begin{matrix}
g_*R^dj_*\mathcal{F} & 0 & 0 \\
\ldots & \ldots & \ldots \\
g_*R^2j_*\mathcal{F} & 0 & 0 \\
g_*R^1j_*\mathcal{F} & 0 & 0 \\
g_*j_*\mathcal{F} & R^1g_*j_*\mathcal{F} & R^2g_*j_*\mathcal{F}
\end{matrix}
$$
We conclude that $R^qf_*\mathcal{F} = g_*R^qj_*\mathcal{F}$ for $q > 2$.
By interlude I we see that the support of $R^qf_*\mathcal{F}$ for $q > 2$
is contained in the set of points of $\mathbf{A}^{d - 1}_K$
whose residue field has transcendence degree $\leq d - q$.
By interlude II
$$
H^p(\mathbf{A}^{d - 1}_K, R^qf_*\mathcal{F}) = 0
\text{ for }p > d - q + \text{cd}(K)\text{ and }q > 2
$$
On the other hand, by Theorem \ref{theorem-higher-direct-images}
we have $R^2f_*\mathcal{F}_{\overline{\eta}} =
H^2(\mathbf{A}^1_{\overline{\eta}}, \mathcal{F}) = 0$
(vanishing by the case of dimension $1$) where $\eta$ is the
generic point of $\mathbf{A}^{d - 1}_K$.
Hence by interlude II again we see
$$
H^p(\mathbf{A}^{d - 1}_K, R^2f_*\mathcal{F}) = 0
\text{ for }p > d - 2 + \text{cd}(K)
$$
Finally, we have
$$
H^p(\mathbf{A}^{d - 1}_K, R^qf_*\mathcal{F}) = 0
\text{ for }p > d - 1 + \text{cd}(K)\text{ and }q = 0, 1
$$
by induction hypothesis. Combining everything we just said
with the Leray spectral sequence
$H^p(\mathbf{A}^{d - 1}_K, R^qf_*\mathcal{F}) \Rightarrow
H^{p + q}(\mathbf{A}^d_K, \mathcal{F})$ we conclude.
\end{proof}

\end{document}
